

\section{Approximability}\label{se:approx}
In this section we introduce various notions of \textit{approximability} that were developed in \cite{ABC26}, with inspiration from a notion appearing in \cite{Turing:approx}, which we first recall.

\subsection{Metric approximability}\label{se:Turing}
We first recall the following definition, due to Turing \cite{Turing:approx}.
\begin{dfn}\label{Turingapprox}
    Let $(G,\circ_G)$ be a group endowed with a left-invariant metric $d$ and let $\varepsilon>0$. We say that $G$ is $\varepsilon$-approximable if there is a \textit{finite subset} $H_\varepsilon \subset G$ endowed with a group structure $\circ_\varepsilon$ such that \begin{enumerate}
        \item for any $x,y\in H_\varepsilon$ we have\footnote{Note the error in the published version of the paper \cite{Turing:approx}.} $$d(x\circ_G y, x\circ_\varepsilon y)<\varepsilon$$ and
        \item for any $x\in G$ there is $C(x)\in H_\varepsilon$ such that $$d(x,C(x))<\varepsilon.$$
    \end{enumerate}
    In the above case we say that $H_\varepsilon$ is an $\varepsilon$-approximation of $G$. We say that $G$ is approximable if it is $\varepsilon$-approximable for every $\varepsilon>0$.
\end{dfn}
\noindent In the same paper \cite{Turing:approx}, Turing proved that the only approximable Lie groups are tori.

Definition \ref{Turingapprox} admits many variations. We will list some possibilities in the following. First, we will not change the basis of the definition (i.e. the fact that it uses \textit{group} structure), and after we will suggest a variation that uses triangulated categories.


We put the above definition into metric context. 
\begin{dfn}
Let $(X,d_X)\subset (Y,d_Y)$ be an inclusion of pseudometric spaces and consider $\varepsilon>0$. A subset $X_\varepsilon\subset Y$ is said to be an $\varepsilon$-net for $X$ in $Y$ if for any $x\in X$ there is $C(x)\in X_\varepsilon$ such that $d_Y(x,C(x))<\varepsilon$. The subset $(X,d_X)\subset (Y,d_Y)$ is said to be totally bounded if for any $\varepsilon>0$ it admits a \textit{finite} $\varepsilon$-net $X_\varepsilon$ in $Y$.
\end{dfn} We recall that a totally bounded metric space has compact completion. Total boundedness is a relaxation of boundedness.

Definition \ref{Turingapprox} says that a metric group is approximable if for any $\varepsilon>0$ it admits a finite $\varepsilon$-net in itself with a very similar group structure. In particular, an approximable metric group is totally bounded.\\
One could replace finiteness of the group $H_\varepsilon$ with finite generation and get a more general notion. Note that, even finite generation is a quite strong notion in this context. For an overview of the study of approability for metric groups see \cite{Thom:approx}. \\

% Composition can be seen as "taking cones" i.e. we might ask that the cones in the category \Tau and in the subcategory (trinagular) generated by \mathcal{F}_\varepsilon are \varepsilon close in the compatible metric.

There are other possible modifications to Definition \ref{Turingapprox}. We will often adopt the following extrinsic viewpoint. \begin{dfn}
    Let $(X,d)$ be a pseudometric space. We say that $(X,d_X)$ is metric $\varepsilon$-approximable in the sense of groups if there is a group $(Y,\circ)$ endowed with a left-invariant metric $d_Y$ such that $(X,d_X)$ isometrically embed in $(Y,d_Y)$ and there is a finite subset $H_\varepsilon\subset Y$ endowed with a group structure such that (1) and (2) in Definition \ref{Turingapprox} are satisfied, with the modification that in (2) the element $x$ is in $X$ and not necessarily in $Y$. If we assume $H_\varepsilon$ to be finitely generated instead of finite, we say that $(D,d)$ is finite-type metric $\varepsilon$-approximable in the sense of groups (in $Y$).

\end{dfn}
In another direction, one could ask the embedding of $(X,d_X$) into $(Y,d_Y)$ to be a quasi-isometry (as we will do in Definition \ref{def:TPCapprox}), or even to assume that the structure of $(X,d_X)$ is that of quasipseudometric (aka Lawvere) space (i.e. of category enriched over the poset $[0,\infty]$, see \cite{Fong_Spivak_2019}).


\subsection{Categorical metric approximability}
The following definition is a version of Definition \ref{Turingapprox} with the concept of \textit{group} replaced by \textit{triangulated category}.
\begin{dfn}[Definition 1.1.1 in \cite{ABC26}] \label{def:approx} Assume that $(X,d_X)\subset (Y,d_Y)$ is an inclusion of metric spaces such that the points in $Y$ are the objects of a triangulated category $\Tau$. For $\varepsilon > 0$, we say that 
$X$ is   {\em categorically metric $\varepsilon$-approximable} in $\Tau$ if there exists a subspace $X_{\varepsilon}\subset Y$ such that the elements of $X_{\varepsilon}$ are the objects of a {\em finitely generated} (in the sense of triangular generation) subcategory of $\Tau$ and, additionally, $X_{\varepsilon}$ is an $\varepsilon$-net for $X$ in $Y$.  If this property is satisfied for all $\varepsilon >0$ we say that $(X,d_X)$ is {\em categorically metric approximable} in $\Tau$. 
\end{dfn}
The following is a weaker version of Definition \ref{def:approx}, and it appears in the Introduction of \cite{ABC26}.
\begin{dfn}\label{def:rapprox}
    Assume that $(X,d_X)\subset (Y,d_Y)$ is an inclusion of metric spaces such that the points in $Y$ are the objects of a triangulated category $\Tau$. For $\varepsilon > 0$, we say that 
$X$ is   {\em retract categorically metric $\varepsilon$-approximable} in $\Tau$ if there exists a subspace $X_{\varepsilon}\subset Y$ such that the elements of $X_{\varepsilon}$ are the objects of a {\em finitely generated} (in the sense of triangular generation) subcategory of $\Tau$ and, additionally, for each $x\in X$
there exists an object $k_{x}$ in $\Tau$, i.e. $k_x\in Y$, such that $x\oplus k_{x}$ satisfies
$d_Y(x\oplus k_{x}, X_{\varepsilon}) < 2\varepsilon$.  If this property is satisfied for all $\varepsilon >0$ we say that $(X,d_X)$ is {\em retract categorically metric approximable} in $\Tau$. 
\end{dfn}

\begin{remark}
    Assume that a pseduometric space $(X,d)$ is categorically metric approximable in a triangulated category $\Tau$. Then $X$ is finite-type metric approximable in the sense of groups in the Grothendieck group $K(\Tau)$ of $\Tau$.
\end{remark}

\begin{remark}
    Definition \ref{Turingapprox} has two flaws. The first one is intrisic in the definition: the triangulated structure on $\Tau$ is completely unrelated to the metric $d_Y$ on its set of objects $Y=\Ob(\Tau)$. We believe this issue in design can be solved by brute force (i.e. assuming that, for instance, $d_Y$ behaves well with respect to taking cones) as well as categorically (i.e. by expressing our structure as some kind of enriched category). \\ Secondly, known examples of categorical metric approximability are limited to the case where $\Tau$ admits a TPC refinement $\msc$ and $d_Y$ is the interleaving metric $\dint$ on $\msc$. In particular, it might seem that at the end of the day Definition \ref{Turingapprox} is not a metric concept, but rather a persistence one. We believe that this is not the case and hope that Definition \ref{Turingapprox} applies in some way to the B-side of the homological mirror symmetry conjecture (i.e. the bounded derived category of coherent sheaves), where no persistence structure seem to exist. Hence the question becomes: how can mirror symmetry be lifted to a quantitative question?
\end{remark}




\subsection{Approximability in a TPC}
In this section we discuss (retract) approximability of a pseudometric space in a TPC. This is a special case of Definition \ref{def:approx}. 

The most basic intrinsic definition is as follows. Recall the operation $\langle\cdot\rangle^\triangle$ defined in Definition \ref{trcompl}.
\begin{dfn}[Definition 2.2.3 in \cite{ABC26}]\label{def:TPCapprox}
    Let $\msc$ be a TPC and consider a subset $\mathcal{L}\subset \Ob(\msc)$. 
    \begin{enumerate}
    \item We say that $\mathcal{L}$ is $\varepsilon$-approximable in $\msc$ if there is a \textit{finite} family $\mathcal{F}_\varepsilon\subset\Ob(\msc)$ such that for any $L\in \mathcal{L}$ we have $$\dint(L,\langle\mathcal{F}_\varepsilon\rangle^\triangle)<\varepsilon$$ 
    We say that $\mathcal{L}$ is approximable in $\msc$ if it is $\varepsilon$-approximable for all $\varepsilon>0$
    \item We say that $\mathcal{L}$ is retract $\varepsilon$-approximable in $\msc$ if there is a \textit{finite} family $\mathcal{F}_\varepsilon\subset\Ob(\msc)$ such that for any $L\in \mathcal{L}$ we have $$\dret(L,\langle\mathcal{F}_\varepsilon\rangle^\triangle)<\varepsilon$$ 
    We say that $\mathcal{L}$ is retract approximable in $\msc$ if it is retract $\varepsilon$-approximable for all $\varepsilon>0$
    \end{enumerate}
\end{dfn}

The first part of the following lemma is immediate, while the second is an exercise in manipulations of triangles.
\begin{lem}
    Let $\msc$ be a TPC. If $\mathcal{L}$ is approximable in $\msc$ then $(\mathcal{L},d_{\text{int}}^\msc\vert_\mathcal{L})$ is categorically metric approximable in the sense of Definition \ref{def:approx}. If $\mathcal{L}$ is retract approximable in $\msc$ then $(\mathcal{L},d_{\text{int}}^\msc\vert_\mathcal{L})$ is categorically metric retract approximable in the sense of Definition \ref{def:rapprox}.
\end{lem}

In the following we will adopt an extrinsic viewpoint on Definition \ref{def:TPCapprox}, meaning that we start with a pseudometric space and (quasi-)isometrically embed it into the objects of a TPC. As it will become apparent when discussing our geometric application, given a space that we want to approximate, we will need the approximating TPC to vary with the approximation accuracy. This has to do with the intrinsic limitations of the machinery of filtered Fukaya categories that will be developed in Chapter \ref{ch:ffuk}, and can be seen as a sort of indeterminacy principle: there is no hope to get finer measurements as what the measuring tool allows. This will be further developed on the algebraic side in \S\ref{s:approximability} and on the geometric side in \S\ref{sysnew}.

\begin{dfn}[Definition 2.2.1 in \cite{ABC26}]\label{def:TPC-approx} A pseudo-metric space $(X,d)$ is TPC-approximable if there exists a constant $A\geq 1$ such that for
  each $\varepsilon >0$ the following structure exists. For each
  $0< \eta <\varepsilon$ there is a TPC, $\msc_{\varepsilon,\eta}$, with two
  properties:
  \begin{itemize}
  \item[i.] There exists an $(A,\eta)$-quasi-isometric embedding (see Remark \ref{rem:def-TPC-approx}):
    $$\Phi_{\varepsilon,\eta}:(X,d) \to (\Ob(\msc_{\varepsilon,\eta}), \sdint^{\ \varepsilon,\eta})$$
  \item[ii.] There exists a {\em finite} family
    $\mathcal{F}_{\varepsilon,\eta} \subset \Ob(\msc_{\varepsilon,\eta})$ with the
    property:
    $$\sdint^{\ \varepsilon,\eta}(\Phi_{\eta}(x), \Ob( \langle\mathcal{F}_{\varepsilon,\eta}
    \rangle ^{\Delta}) ) < \varepsilon, \ \forall x\in X$$
  \end{itemize} 
  where $\sdint^{\ \varepsilon,\eta}$ is the shift invariant interleaving pseudo-metric associated to
  $\msc_{\varepsilon,\eta}$ introduced in \S\ref{dint}.
\end{dfn}

For fixed
$\varepsilon$, the data $(\Phi, \F)$ with  $\Phi= \{\Phi_{\varepsilon,\eta}\}_{0<\eta <\varepsilon}$,
$\F=\{\F_{\varepsilon,\eta} \}_{0<\eta<\varepsilon}$ is called 
TPC $\varepsilon$-{\em approximating data} for $(X,d)$. Thus, $(X,d)$ is
TPC-approximable if $\varepsilon$-TPC-approximating data for $(X,d)$
exists for each $\varepsilon>0$. It can certainly happen that multiple
choices of such data exist for the same $(X,d)$. When one of the two parameters $\varepsilon$ or $\eta$ are clear from the context we omit it from the notation. For instance, for fixed $\varepsilon$, we write $\Phi_{\eta}$ for the relevant quasi-isometric embeddings , $\sdint^{\ \eta}$ for the respective pseudo-metrics and so forth. 


\

This definition  is the TPC  version of Definition \ref{def:approx}. 
We will also  use in the paper the notion of  {\em TPC retract approximability}.  TPC retract approximability is defined just as above
without any modifications for point i. above but by using $ \sdret$ at
point ii.  (we emphasize that the {\em only} change is at point ii. in
Definition \ref{def:TPC-approx}). Pairs $(\Phi, \F)$ as above but ensuring retract approximability (in other words, using $\sdret$ instead of $\sdint$ for condition ii in the definition) will be called TPC {\em retract $\varepsilon$-approximating data}.

\

Point c. in Remark \ref{rem:relation-TPCapp} below explains more precisely how  TPC  approximability, as just defined, fits with the notion of categorical metric approximability introduced in Definition \ref{def:approx}.


\begin{rem}\label{rem:relation-TPCapp}
  a. Recall that an $(A,B)$-quasi-isometric embedding
  $$(X,d) \to (Y,d')$$ between two metric spaces is a
  map $\Phi : X\to Y$ such that:
  $$\frac{1}{A}\  d(x,y)-B \leq d'(\Phi(x),\Phi(y))\leq A\ d(x,y)+B~.~$$
  Thus, the first point of the Definition \ref{def:TPC-approx} shows that, when $\varepsilon$ is
  fixed, the metrics
  $\sdint^{\ \eta}$ defined on the objects of $\msc_{\eta}$, when
  restricted to $X$, get closer and closer to a (pseudo)-metric
  equivalent to $d$, when $\eta\to 0$. More precisely, we can put for
  each $x,y\in X$
  $$\hat{d}(x,y)=\limsup_{\eta\to 0} \sdint^{\ \eta} (x,y)~.~$$
  This gives a pseudo-metric defined on $X$ and point i. of the
  Definition implies that this $\hat{d}$ is equivalent to $d$, the
  original metric on $X$.

  b. A $(A,\eta)$-quasi-isometric embedding is also a
  $(A,\eta')$-quasi-isometric embedding whenever $\eta'\geq \eta$. As
  a result, to verify the property in the definition, one only needs
  to show the existence of $\msc_{\eta}$, $\Phi_{\eta}$, and of
  the families $\F_{\varepsilon} =\{\F_{\varepsilon,\eta}\}$ whenever $\eta$ is small enough.


  c. In Definition \ref{def:TPC-approx} we have not imposed any
  relation tying the categories $\msc_{\varepsilon,\eta}$, when
  $\eta \to 0$, nor did we assume any relation among the respective
  families $\F_{\varepsilon,\eta}$.  In practice, our constructions in
  the case of Fukaya categories produce categories that satisfy a
  property more precise than the one in Definition
  \ref{def:TPC-approx}, in the sense that, for fixed $\varepsilon >0$, the
  categories $\msc_{\varepsilon,\eta}$, with $\eta\to 0$, are related by
  certain comparison functors and the corresponding families
  $\F_{\varepsilon,\eta}$ correspond one to the other through these
  functors. This more refined property is more technical to state and
  is formulated in Definition \ref{d:approx-sys}. Further below we
  will also discuss the dependence of our constructions relative to $\varepsilon$.
  
\end{rem}

\subsection{Approximability in a system of TPCs} \label{s:approximability}
%\subsubsection{Approximability within a system} \label{sb:approx-sys}
This is Section 2.5.10 in \cite{ABC26}. Let $\widehat{\msc}$ be a system of TPC's with increasing accuracy, as
defined in \S\ref{s:sys-tpc}. We will use here the equivalence relation
  $\sim_{\widehat{\msc}}$ on
  $\Ob^{\text{tot}}(\widehat{\msc}) = \coprod_{p\in \mathcal{P}}
  \Ob(\msc_p)$ as defined in~\S\ref{sb:sys-inv} on
  page~\pageref{p:equiv-ob} and the notation from that section.
\begin{dfn} \label{d:approx-sys} Let
  $\mathcal{L} \subset \widetilde{\Ob}(\widehat{\msc})$ and let
  $\mathcal{F}_1, \ldots, \mathcal{F}_l \in
  \widetilde{\Ob}(\widehat{\msc})$ be a finite collection (of equivalence
  classes of objects), and let $\varepsilon>0$. We say that $\mathcal{L}$
  is $\varepsilon$-approximable by
  $\{\mathcal{F}_1, \ldots, \mathcal{F}_l\}$ if there exists
  $\delta>0$ and $c_{\varepsilon}>0$ such that for every
  $p \in \mathcal{P}$ with $\nu(p) < \delta$ we have
  $${d}_{\text{int}} \Bigl(\mathcal{L}_p,
  \Ob \langle \{F_1, \ldots, F_l\}_p \rangle^{\Delta} \Bigr)
  < \varepsilon + c_{\varepsilon}\nu(p).$$
  Let $\mathcal{F} \subset \widetilde{\Ob}(\widehat{\msc})$. We say that
  $\mathcal{L}$ is approximable by $\mathcal{F}$ if for every
  $\varepsilon>0$ there exists a finite collection
  $\mathcal{F}_1, \ldots, \mathcal{F}_l \in \mathcal{F}$ which
  $\varepsilon$-approximates $\mathcal{L}$ and the constants $c_{\varepsilon}$ 
  are bounded, independently of $\varepsilon$.

  \end{dfn}
  Recall that 
  %$\{F_1, \ldots, F_l\}_p^{\Sigma}$ stands
  %for completing the set of objects
  %$\{F_1, \ldots, F_l\}_p \subset \Ob(\msc_p)$ with respect to all
  %shifts, and d
  $\langle \mathcal{O} \rangle^{\Delta}$ stands for the
  smallest sub-TPC of $\msc_p$ containing the set of
  objects $\mathcal{O}$. Sometimes we will use the wording {\em
    $\{\mathcal{F}_1, \ldots, \mathcal{F}_l\}$ $\varepsilon$-approximates
    $\mathcal{L}$} and {\em $\mathcal{F}$ approximates $\mathcal{L}$}.

\

Similarly we also define retract approximability as follows. 
\begin{dfn} \label{d:approx-sys-ret} Let $\widehat{\msc}$,
  $\mathcal{L}$, $\mathcal{F}_1, \ldots, \mathcal{F}_l$ be as in
  Definition~\ref{d:approx-sys} and let $\varepsilon>0$. We say that
  $\mathcal{L}$ is $\varepsilon$-retract-approximable by
  $\{\mathcal{F}_1, \ldots, \mathcal{F}_l\}$ if there exists
  $\delta>0$, $c_{\varepsilon}>0$, such that for every $p \in \mathcal{P}$ with
  $\nu(p) < \delta$ we have
  $${d}_{\text{r-int}} \Bigl(\mathcal{L}_p,
  \Ob \langle \{F_1, \ldots, F_l\}_p \rangle^{\Delta} \Bigr)
  < \varepsilon +c_{\varepsilon}\nu(p).$$

  Let $\mathcal{F} \subset \widetilde{\Ob}(\widehat{\msc})$. We say
  that $\mathcal{L}$ is retract-approximable by $\mathcal{F}$ if for
  every $\varepsilon>0$ there exists a finite collection
  $\mathcal{F}_1, \ldots, \mathcal{F}_l \in \mathcal{F}$ which
  $\varepsilon$-retract-approximates $\mathcal{L}$ and the constants
  $c_{\varepsilon}$ are uniformly bounded.
\end{dfn}

% \begin{rem}\label{rem:different_approx} It is useful at this point to
%   relate our two notions of approximability, in Definitions
%   \ref{def:TPC-approx} and \ref{d:approx-sys}.  In our geometric
%   applications, the set $\mathcal{L}$ will be a set of Lagrangian
%   submanifolds $\mathcal{L}ag(M)$ and the category $\msc_{p}$ will be
%   a certain TPC associated with a filtered Fukaya category defined using
%   perturbation data that is identified by the parameter $p$.  The
%   inclusion $\mathcal{L}ag(M)\subset \msc_{p}$ is the Yoneda
%   embedding. The set $\mathcal{L}ag(M)$ is endowed with the spectral
%   metric $d_{\gamma}$ (that will be recalled later) and this metric is
%   closely related to the stabilized interleaving metric $\sdint$ from
%   \eqref{eq:sdint} restricted to $\mathcal{L}ag(M)$.  In view of this,
%   we will see that $\varepsilon$ - approximability in the sense of
%   Definition \ref{d:approx-sys} implies $\varepsilon'$-approximability
%   for all $\varepsilon'>\varepsilon$ in the sense of Definition
%   \ref{def:TPC-approx}. The reason for this discrepancy in the
%   parameters $\varepsilon,\varepsilon'$ is related to the presence of the
%   term $c_{\varepsilon}\nu(p)$ in Definition \ref{d:approx-sys}.  Given
%   that these properties are supposed to be satisfied for all
%   $\varepsilon,\varepsilon' >0$ and $\nu(p)$ can be taken as small as
%   desired, this discrepancy is essentially irrelevant. However, for a
%   fixed $\varepsilon$, including the term $c_{\varepsilon} \nu(p)$ in the
%   upper bound in Definition \ref{d:approx-sys} is useful for questions
%   having to do with minimizing the cardinality of the set
%   $\{\mathcal{F}_{i}\}$ for that fixed $\varepsilon$.  A similar remark
%   applies to retract-approximability.
% \end{rem}

\subsection{Stabilizing functors} \label{sbsb:sfunc} This is Section 2.5.11 in \cite{ABC26}. Let
  $\widehat{\msc} = \{ \msc_p\}_{p \in \mathcal{P}}$ be a system of
  PC's with increasing accuracy and let
  $\widehat{\msf}: \widehat{\msc} \longrightarrow \widehat{\msc}$ be
  an endo-functor of such systems (see~\S\ref{sbsb:func-sys-pc}). Let
  $\mathscr{B} \subset \Ob^{\text{tot}}(\widehat{\msc})$. We say that
  $\widehat{\msf}$ is {\em $\mathscr{B}$-stabilizing} if the following
  holds for every $p \in \mathcal{P}$:
  \begin{enumerate}
  \item \label{i:stab-1} For every $A \in \mathscr{B}_p$ we have
    $(A,p) \sim_{\widehat{\msc}} (\msf_p(A), \widehat{\msf}(p))$,
    where $\sim_{\widehat{\msc}}$ is the equivalence relation
    from~\S\ref{sb:sys-inv}, page~\pageref{p:equiv-ob}.
  \item \label{i:stab-2} For every $A', A'' \in \mathscr{B}_p$ there
    exists $q \in \mathcal{P}$ with $p, \widehat{\msf}(p) \preceq q$
    and two $0$-isomorphisms
    $\sigma_{A'}: \msh_{\widehat{\msf}(p), q} \msf_p A'
    \longrightarrow \msh_{p,q}A'$,
    $\sigma_{A''}: \msh_{\widehat{\msf}(p), q} \msf_p A''
    \longrightarrow \msh_{p,q}A''$ such that for every
    $u \in \hom_{\msc_p}(A',A'')$ we have:
    \begin{equation} \label{eq:stab-2} \sigma_{A''} \circ
      \msh_{\widehat{\msf}(p), q} \circ \msf_p(u) = \msh_{p,q}(u)
      \circ \sigma_{A'}.
    \end{equation}
  \end{enumerate}

\begin{rem} \label{r:stab-func} 
Condition~\eqref{i:stab-1}
    above is equivalent to requiring that the map
    $\widetilde{\msf}: \widetilde{\Ob}(\widehat{\msc}) \longrightarrow
    \widetilde{\Ob}(\widehat{\msc})$ (see Remark~\ref{r:d-limit})
    stabilizes $\mathscr{B}$ (i.e.~$\widetilde{\msf}$ sends
    $\mathscr{B}$ to itself and restricts to the identity map
    $\mathscr{B} \longrightarrow \mathscr{B}$).

  Condition~\eqref{i:stab-2} can be rephrased as follows. By a
    slight abuse of notation denote for every $p \in \mathcal{P}$ by
    $\mathscr{B}_p \subset \msc_p$ the full persistence subcategory
    with objects $\mathscr{B}_p$. We require that given
    $p \in \mathcal{P}$, for large enough (with respect to $\preceq$)
    $q \in \mathcal{P}$ the two restricted PC functors
    $$\msh_{\widehat{\msf}(p), q} \circ \msf_p|_{\mathscr{B}_p}, \, 
    \msh_{p,q}|_{\mathscr{B}_p}: \mathscr{B}_p \longrightarrow
    \mathscr{B}_q$$ are naturally $0$-isomorphic in the category of
    persistence functors
    $\mathscr{B}_p \longrightarrow \mathscr{B}_q$.
  \end{rem}

  In the presence of approximating families in TPC's,
    stabilizing functors have a remarkable metric property.
    \begin{prop} \label{p:stab-func} Let $\widehat{\msc}$ be a system
      of TPC's, $\mathcal{L} \subset \widetilde{\Ob}(\widehat{\msc})$
      and
      $\mathcal{F}_1, \ldots, \mathcal{F}_l \in
      \widetilde{\Ob}(\widehat{\msc})$ an $\varepsilon$-approximating
      family for $\mathcal{L}$, as in Definition~\ref{d:approx-sys}.
      Let
      $\widehat{\Phi}: \widehat{\msc} \longrightarrow \widehat{\msc}$
      be a TPC endo-functor of systems and assume that
      $\widehat{\Phi}$ is
      $\{\mathcal{F}_1, \ldots, \mathcal{F}_l\}$-stabilizing. Then for
      every $\widetilde{L} \in \mathcal{L}$ we have
      $\widetilde{d}_{\text{int}}(\widetilde{L},
      \widetilde{\Phi}(\widetilde{L})) \leq \varepsilon$.
  \end{prop}